% from ega4-14-fr.tex
\oldpage[IV-3]{218}

  \item $f$ est universellement ouvert dans un voisinage de $f^{-1}(y)$.
  \item $f$ est ouvert dans un voisinage de $f^{-1}(y)$.
  \item Toute composante irréductible de $X$ rencontrant $f^{-1}(y)$ domine $Y$.
  \item Pour tout point $x\in f^{-1}(y)$, fermé dans $f^{-1}(y)$, et toute
  composante irréductible $X_i$ de $X$ contenant $x$, il existe une partie
  irréductible $X'$ de $X$, localement fermée dans $X$, contenant $x$, contenue
  dans $X_i$, et telle que la restriction $X'\to Y$ de $f$ soit un morphisme
  quasi-fini et dominant.
\end{enumerate}
\end{proposition}

L'équivalence de $a)$, $b)$, $c)$ et $d)$ a déjà été démontrée (14.3.8). Il est
clair que $e)$ entraîne $d)$, car pour toute composante irréductible $X_i$ de $X$
rencontrant $f^{-1}(y)$, il existe dans $X_i\cap f^{-1}(y)$ un point fermé dans cet
espace (5.1.11). Enfin, pour prouver que $c)$ entraîne $e)$, on peut se borner au
cas où $f$ est un morphisme de type fini~; considérons un point fermé $x$ de
$f^{-1}(y)$ et soit $X_i$ une composante irréductible de $X$ contenant $x$. Comme
la restriction $f_i:X_i\to Y$ de $f$ à $X_i$ est un morphisme dominant, il résulte
de l'équivalence de $c)$ et $d)$ pour $f_i$ que ce morphisme est ouvert aux points
génériques de $X_i\cap f^{-1}(y)$. Il résulte donc de (14.2.2) que $f_i$ est
équidimensionnel au point $x$, et on conclut alors à l'aide de (14.5.3).

\begin{remark}[14.5.7]
\EGAReaderTarget{ega4.c04.rem.14-5-7}{IV.14.5.7-fr}
Si, dans l'énoncé de (14.5.5), on suppose que $Y=\operatorname{Spec}(A)$, où $A$ est un anneau
de valuation discrète complet, et que $y$ est le point fermé de $Y$, on peut en
outre supposer que $X'=\operatorname{Spec}(A')$, où $A'$ est un anneau de valuation discrète qui
est une $A$-algèbre finie, comme le montre la démonstration de (14.5.4) et le fait
qu'un anneau local intègre régulier et de dimension 1 est un anneau de valuation
discrète (\textbf{II}, 7.1.6).
\end{remark}

\begin{proposition}[14.5.8]
\label{IV.14.5.8-fr}
Soient $Y$ un préschéma localement noethérien, $f:X\to Y$ un morphisme localement
de type fini, $y$ un point de $Y$.

Pour tout $Y$-préschéma $Y'$, posons $X(Y')=\operatorname{Hom}_Y(Y',X)$. Soit $x$ un point de
$f^{-1}(y)$~; afin que $f$ soit universellement ouvert au point $x$, il faut et il
suffit que la condition suivante soit vérifiée~:

Pour tout anneau de valuation discrète complet $A$, de corps résiduel
algébriquement clos, tout morphisme $g:Y'=\operatorname{Spec}(A)\to Y$ tel que l'image par $g$ du
point fermé $y'$ de $Y'$ soit égal à $y$, et tout élément
$u_0\in X(\operatorname{Spec}(k(y')))$ tel que $u_0(y')=x$, il existe un anneau de valuation
discrète $B$, un homomorphisme local $A\to B$ faisant de $B$ une $A$-algèbre finie,
et, si l'on pose $Z'=\operatorname{Spec}(B)$, un élément $u\in X(Z')$ tels que, si $z'$ est le
point fermé de $Z'$, le diagramme
\[
  \xymatrix{
    \operatorname{Spec}(k(z')) \ar[r] \ar[d]^{\ell} & Z' \ar[d]^{u} \\
    \operatorname{Spec}(k(y')) \ar[r]_{u_0} & X
  }
\]
soit commutatif.
\end{proposition}

